\section{Conclusiones}\label{sec:conclusiones}

\begin{itemize}
  \item El reto es de interpolación. Hay que estimar el rendimiento de 59 parcelas a partir de 138 vecinas del mismo
        año, y la división entre ambos conjuntos se hizo al azar.
  \item Las imágenes satelitales son abundantes y cubren todo 2025. Lo escaso son los rendimientos, y no existe una
        fuente pública que los complemente.
  \item Los rendimientos parecen venir de dos formas de medir, separadas por municipio y no por estado. Si se
        confirma, esa variable explicará buena parte de las diferencias entre parcelas.
  \item Los rendimientos los declaró cada productor, sin validación oficial. Parte del error es ruido de etiqueta, probablemente sobre todo en el régimen redondeado y en las parcelas pequeñas, y ese ruido pone un piso a lo que se puede lograr.
  \item Unos 17 de los 23 índices del conjunto Básico son independientes. Los demás son transformaciones de otro y conviene conservar sólo uno de cada pareja.
  \item En 2025 importó más cuándo llovió que cuánto, y sus efectos fueron de parcela y no de zona.
  \item Con 138 parcelas, las mejoras pequeñas pueden ser azar. Toda idea pasa por la misma regla de decisión.
  \item Lo observado hasta ahora es preliminar y debe comprobarse de forma independiente antes de apoyar una
        decisión.
\end{itemize}
